Il patrimonio librario posseduto da utenti privati rappresenta una risorsa culturale distribuita, spesso significativa ma difficilmente accessibile al di fuori della sfera personale del proprietario. In assenza di strumenti di catalogazione e pubblicazione, una collezione privata rimane non individuabile dall'esterno, e la distanza tra disponibilità fisica e accessibilità effettiva impedisce la scoperta di volumi potenzialmente utili a una comunità di lettori.

Questa tesi affronta il problema progettando e realizzando un prototipo web geolocalizzato per rendere catalogabili, individuabili e condivisibili patrimoni librari appartenenti a utenti privati. Il sistema consente di registrarsi, pubblicare libri con metadati bibliografici, renderli ricercabili tramite criteri testuali e geografici, visualizzarli su mappa interattiva e ricevere richieste di consultazione, prestito o informazione. Un aspetto centrale della soluzione riguarda la tutela della posizione: il progetto distingue tra coordinate precise, usate internamente, e coordinate pubbliche approssimate, utilizzate per la discovery territoriale, in modo da conciliare prossimità geografica e minimizzazione dell'esposizione dei dati sensibili.

L'architettura è basata su un monorepo modulare con Next.js, PostgreSQL/PostGIS, Prisma, Auth.js e MapLibre GL JS. La validazione ha combinato controlli statici, build di produzione e scenari funzionali manuali, confermando la correttezza dei flussi principali: autenticazione, catalogazione, geolocalizzazione, ricerca per prossimità e richieste tra utenti. I risultati dimostrano come un requisito di localizzazione possa essere integrato in una piattaforma di catalogazione senza esporre direttamente la posizione precisa degli utenti, fornendo un modello realizzabile e verificabile di discovery territoriale per risorse fisiche distribuite.
